\chapter{机械波}\label{ch:g12:mechanical-waves}

周六傍晚，满座的体育场里，一股“人浪”
不到一分钟便绕过看台，可是没有人离开自己的座位。
石子激起的圆环穿过池面，而软木塞只是上下浮动。
本章研究这个“旅行者”本身——波——以及驯服它的四个量：
速率、延迟、周期、波长——这些量足以让人
从一张纸带上读出地震。

\section{扰动向前传播，介质留在原地}

\begin{definition}[机械波]\label{def:g12:mechanical-waves:wave}
称\emph{机械波}\index{机械波}为一种扰动，
它在实体介质——绳子、水、空气、岩石——中传播，
传递能量而\emph{不迁移物质}：介质的每一小部分只在原地短暂振动，
将运动传给相邻部分，然后回到原位。
\end{definition}

\begin{example}[物质不随波迁移]\label{ex:g12:mechanical-waves:transport}
在体育场的人浪中，“介质”是观众：传播的是
起立的动作，而每位观众都留在自己的座位上。
在池塘中，圆环向外扩展，而漂浮的软木塞只在原地上下浮动。
\end{example}

\begin{definition}[横波与纵波]\label{def:g12:mechanical-waves:transverse}
波称为\emph{横波}\index{横波}，即介质的运动与
传播方向\emph{垂直}：一个凸起沿被抖动的绳子快速传播，
而绳上每一点只上下运动。波称为
\emph{纵波}\index{纵波}，即介质的运动\emph{沿}
传播方向：一段被压缩的线圈沿被推动的弹簧向下传播，
每个线圈只在原地振动。
\end{definition}

\begin{omfigure}
\begin{tikzpicture}[font=\small]
  \node[left] at (-0.15,2.3) {横波};
  \draw[omDef, thick, smooth, samples=100, domain=0:8.4]
    plot (\x, {2.3 + 0.42*sin(130*\x)});
  \draw[omThm, Stealth-Stealth] (6.92,1.78) -- (6.92,2.82)
    node[right, yshift=-14pt] {运动};
  \draw[-Stealth, thick] (3.4,3.15) -- (5.0,3.15)
    node[midway, above] {传播};
  \node[left] at (-0.15,0) {纵波};
  \foreach \i in {0,...,41}
    \draw[omDef, thick]
      ({0.2*\i + 0.26*sin(32*\i)},0.32) -- ({0.2*\i + 0.26*sin(32*\i)},-0.32);
  \node[below] at (3.4,-0.35) {密部};
  \node[below] at (6.65,-0.35) {疏部};
  \draw[omThm, Stealth-Stealth] (0.75,0.55) -- (1.55,0.55)
    node[midway, above] {运动};
\end{tikzpicture}

{\small 横波（绳子）：介质的运动与传播方向垂直。
纵波（弹簧）：密部传播，线圈只在原地振动。}
\end{omfigure}

\begin{remark}[声音是纵波]\label{rem:g12:mechanical-waves:sound}
声音是典型的纵波：被压缩的空气——
压强较高（见\cref{ch:g10:pressure}）——从声源传到耳膜；
\cref{ch:g12:sound-acoustics}专门讨论这一内容。地震会同时产生两种波：
P 波压缩岩石，S 波使岩石横向晃动。
\end{remark}

\section{速率与延迟}

\begin{definition}[传播速度]\label{def:g12:mechanical-waves:speed}
波的\emph{传播速度}\index{传播速度}是扰动经过的
距离与所用时间之比，
即 $v = d/t$。它是行进的\emph{波形}的速率，
而非介质的速率。
\end{definition}

\begin{omfigure}
\begin{tikzpicture}[font=\small]
  \draw[-Stealth] (-0.4,0) -- (8.9,0) node[below] {$x$ (\unit{m})};
  \foreach \x in {0,1,...,8}
    \draw (\x,0.05) -- (\x,-0.05) node[below] {$\x$};
  \draw[omDef, thick, smooth, samples=90, domain=-0.2:8.6]
    plot (\x, {1.1*exp(-3.5*(\x-1.0)^2)});
  \draw[omThm, thick, dashed, smooth, samples=90, domain=-0.2:8.6]
    plot (\x, {1.1*exp(-3.5*(\x-3.4)^2)});
  \node[omDef, above] at (1.0,1.15) {$t = 0$};
  \node[omThm, above] at (3.75,1.15) {$t = \qty{0.30}{s}$};
  \draw[omProp, Stealth-Stealth] (1.0,1.75) -- (3.4,1.75)
    node[midway, above] {$d = v\,\Delta t$};
\end{tikzpicture}

{\small 一个绳波脉冲拍摄了两张照片：波峰移动了 \qty{2.4}{m}
用时 \qty{0.30}{s}，故 $v = \qty{8.0}{m/s}$——波形在传播，
而绳子保持原位。}
\end{omfigure}

\begin{proposition}[介质决定速率]\label{prop:g12:mechanical-waves:medium}
对于小扰动，传播速度仅取决于
\emph{介质}，而与扰动的形状或大小无关。在
绷紧的绳子上，$v$ 随张力增大而增大，
随每米质量增大而减小：绷紧而轻则快，松弛而重则慢。\admitted
\end{proposition}

\begin{remark}[公式的出处]\label{rem:g12:mechanical-waves:ropeformula}
绳中波速的精确公式在本科一年级分册中推导；
本学年，上述实验事实已经足够。
\end{remark}

\begin{example}[声音的速率]\label{ex:g12:mechanical-waves:speeds}
声音在空气中的速率为 \qty{340}{m/s}，在水中为 \qty{1500}{m/s}，
在钢中为 \qty{5900}{m/s}——介质越硬，速率越大。由此可得西部电影里的经典镜头：
把耳朵贴在铁轨上，钢轨会比空气更早传来火车的声音。
\end{example}

\begin{definition}[延迟]\label{def:g12:mechanical-waves:delay}
速率为 $v$ 的波传到距离为 $d$ 的一点时所经历的
\emph{延迟}\index{延迟} $\tau = d/v$：相距 $d$ 的两点
重复相同的运动，较远的一点滞后 $\tau$——
恰是\cref{ch:g10:signals-and-waves}中回声测距的逆用。
\end{definition}

\begin{proposition}[无相互作用的叠加]\label{prop:g12:mechanical-waves:superposition}
两列波相遇重叠时，介质的位移在每一时刻
都等于各列波单独存在时所引起位移的\emph{和}；
相遇之后，各列波仍保持原状——
这一结论在本科二年级分册中由运动方程的线性导出。\admitted
\end{proposition}

\begin{example}[相遇的脉冲]\label{ex:g12:mechanical-waves:crossing}
从绳子两端各送出一个向上的凸起：两者相遇时，绳子
短暂升高到两者高度之和；随后，两个完整的凸起
重新分开，各自穿出继续前进——正如两段交错的谈话各自传到听众耳中。
\end{example}

\section{周期波：双重周期性}

\begin{definition}[周期波]\label{def:g12:mechanical-waves:periodic}
当波源以周期 $T$ 重复其运动，频率
为 $f = 1/T$（见\cref{ch:g10:signals-and-waves}），则波源向介质
输入一列\emph{周期波}\index{周期波}：每个质点都以周期 $T$ 重复自身
的运动，且各质点都比相邻质点滞后一个传播延迟。
\end{definition}

\begin{definition}[波长]\label{def:g12:mechanical-waves:wavelength}
周期波的一个快照呈现在空间中重复的图样；其
空间周期——即相邻波峰之间的距离——就是
\emph{波长}\index{波长} $\lambda$，单位为米。周期波
具有\emph{双重周期性}\index{双重周期性}：在各点处，时间上的周期为 $T$，
在每一时刻，空间上的周期为 $\lambda$。
\end{definition}

\begin{theorem}[波的基本关系式]\label{thm:g12:mechanical-waves:relation}
速率为 $v$、周期为 $T$、频率为 $f$ 的周期波，其
波长 $\lambda = v\,T = v/f$。
\end{theorem}

\begin{proof}
在一个周期内，图样向前移动 $v\,T$；但波源
——从而每个质点——此时回到起始状态，因此
移动后的图样与原图样重合：重复距离为 $v\,T$。
\end{proof}

\begin{omfigure}
\begin{tikzpicture}[font=\small]
  \draw[-Stealth] (-0.25,0) -- (4.35,0) node[below] {$x$};
  \draw[-Stealth] (0,-1.1) -- (0,1.55) node[left] {$y$};
  \draw[omDef, thick, smooth, samples=110, domain=0:4.1]
    plot (\x, {0.8*sin(160*\x)});
  \draw[omProp, Stealth-Stealth] (0.5625,1.05) -- (2.8125,1.05)
    node[midway, above] {$\lambda$};
  \node at (2.05,-1.55) {快照 $y(x)$，$t$ 固定};
  \begin{scope}[xshift=5.6cm]
    \draw[-Stealth] (-0.25,0) -- (4.35,0) node[below] {$t$};
    \draw[-Stealth] (0,-1.1) -- (0,1.55) node[left] {$y$};
    \draw[omThm, thick, smooth, samples=110, domain=0:4.1]
      plot (\x, {0.8*sin(160*\x)});
    \draw[omProp, Stealth-Stealth] (0.5625,1.05) -- (2.8125,1.05)
      node[midway, above] {$T$};
    \node at (2.05,-1.55) {信号 $y(t)$，$x$ 固定};
  \end{scope}
\end{tikzpicture}

{\small 同一条正弦曲线的两幅图：快照在空间中每隔 $\lambda$ 重复，
单点信号在时间上每隔 $T$ 重复一次；$\lambda = v\,T$。}
\end{omfigure}

\begin{remark}[频率由波源决定]\label{rem:g12:mechanical-waves:source}
波从一种介质进入另一种介质——例如声音由空气进入水中——
频率保持不变（由波源决定），速率
则发生改变；由 $\lambda = v/f$ 可知波长随之改变。
\end{remark}

\begin{example}[标准音 A：空气中与水中]\label{ex:g12:mechanical-waves:concertA}
在\cref{ch:g10:signals-and-waves}中的 \qty{440}{Hz} 标准音 A，
其波长在空气中为 $340/440 = \qty{0.77}{m}$，在水中为 $1500/440 =
\qty{3.4}{m}$。同一个音，被速率更大的介质拉长。
\end{example}

\begin{definition}[同相与反相]\label{def:g12:mechanical-waves:phase}
相距整数个波长（$d = k\lambda$）的两点，
在任一时刻的运动完全相同：它们
处于\emph{同相}\index{同相}。相距
$d = (k + \tfrac12)\lambda$ 的两点运动始终相反——波峰对
波谷：它们处于\emph{反相}\index{反相}。
\end{definition}

\begin{method}[读取双重周期性]\label{met:g12:mechanical-waves:reading}
\begin{enumerate}
  \item 对固定点处的记录 $y(t)$，由相邻波峰读出 $T$；
        $f = 1/T$。
  \item 对某一时刻的快照 $y(x)$，读出 $\lambda$
        （相邻波峰之间的距离）。
  \item 推得 $v = \lambda/T = \lambda f$；
        若有直接测时，则用 $v = d/\tau$ 核对。
  \item 切勿混淆两种图像：一条坐标轴以秒为单位，
        另一条以米为单位。
\end{enumerate}
\end{method}

\section{波前与衰减}

\begin{definition}[波前]\label{def:g12:mechanical-waves:wavefront}
\emph{波前}\index{波前}是一条线（或一个面），其上各点
在同一时刻被波到达——例如池塘水面上的一条波峰线。
由点波源发出的波前是向外扩展的圆；远离
任何波源的地方，波前近乎直线。波沿
与波前垂直的方向传播，相邻波峰相距 $\lambda$。
\end{definition}

\begin{definition}[衰减]\label{def:g12:mechanical-waves:attenuation}
波在传播过程中振幅的减小，称为其
\emph{衰减}\index{衰减}。两个原因共同作用：一是
圆形波前的能量分布在越来越长的波峰上，二是
介质在波通过时吸收少量能量。直线波前可以避开
第一个原因——这正是远方风暴的涌浪能横跨大洋的缘故。
\end{definition}

\begin{omfigure}
\begin{tikzpicture}[font=\small]
  \fill (0,0) circle (2pt);
  \foreach \r/\c in {0.5/100, 1.0/82, 1.5/64, 2.0/48, 2.5/34}
    \draw[thick, omDef!\c!white] (0,0) circle (\r);
  \draw[omProp, Stealth-Stealth] (315:1.5) -- (315:2.0)
    node[midway, below left, inner sep=1pt] {$\lambda$};
  \draw[-Stealth, omThm, thick] (25:2.55) -- (25:3.25)
    node[right] {$v$};
  \foreach \x/\c in {4.9/100, 5.5/90, 6.1/80, 6.7/70}
    \draw[thick, omDef!\c!white] (\x,-1.5) -- (\x,1.5);
  \draw[-Stealth, omThm, thick] (6.9,0) -- (7.6,0) node[right] {$v$};
\end{tikzpicture}

{\small 圆形波前（相邻两条相距一个波长）随能量
在更长的波峰上分散而逐渐减弱；直线波前（右）减弱较慢。}
\end{omfigure}

\section{练习}

\begin{exercise}[$\star$]\label{exo:g12:mechanical-waves:1}
体育场人浪沿长 \qty{320}{m} 的环形看台传播，绕行一周用时 \qty{40}{s}。
(a) 求其速率。(b) 是否有观众随之绕行？
若不是，那么绕行的是什么？(c) 每位观众起立的时间比
距其 \qty{12}{m} 远的邻座晚 \qty{1.5}{s}，请检验这一说法是否自洽。
\end{exercise}

\begin{exercise}[$\star$]\label{exo:g12:mechanical-waves:2}
横波还是纵波？(a) 上下抖动绳子产生的脉冲；
(b) 沿弹簧传递的压缩；(c) 空气中的声音；(d) 池塘
中的涟漪（水面竖直起伏，波纹水平传播）；(e)
地震 P 波（沿传播方向压缩）和 S 波（左右抖动）。
\end{exercise}

\begin{exercise}[$\star$]\label{exo:g12:mechanical-waves:3}
闪电之后 \qty{4.0}{s} 你听到雷声。(a) 雷击处
距你多远（$v = \qty{340}{m/s}$）？(b) 一名潜水员
在水下听到的闪光与爆响之间的延迟与此相同（$v = \qty{1500}{m/s}$），
则声源距离多远？(c) 把延迟换算为距离前，须知道什么？
\end{exercise}

\begin{exercise}[$\star$]\label{exo:g12:mechanical-waves:4}
求频率为 \qty{440}{Hz} 的声波在空气和水中的波长。
声波由空气进入水中时，$f$、$v$
、$\lambda$ 中哪些改变，哪些保持不变？
\end{exercise}

\begin{exercise}[$\star$]\label{exo:g12:mechanical-waves:5}
池塘中相邻两个波峰相距 \qty{12}{cm}，波环
以 \qty{0.30}{m/s} 的速率扩展。求该波的波长、周期
和频率。
\end{exercise}

\begin{exercise}[$\star\star$]\label{exo:g12:mechanical-waves:6}
一条长绳上的脉冲两次被拍照，两次拍照相隔 \qty{0.30}{s}，
波峰先位于 $x = \qty{1.0}{m}$，后位于 \qty{3.4}{m} 处。(a) 求其速率。
(b) 脉冲何时到达 $x = \qty{7.0}{m}$ 处的墙？(c) 一条
在 $x = \qty{5.0}{m}$ 处打结的彩带，脉冲经过时它如何运动？
\end{exercise}

\begin{exercise}[$\star\star$]\label{exo:g12:mechanical-waves:7}
两个向上的脉冲，高度分别为 \qty{3.0}{cm} 和 \qty{2.0}{cm}，相向
各以 \qty{2.0}{m/s} 的速率传播，两波峰相距 \qty{4.0}{m}。
(a) 两波峰何时重合？(b) 此时绳子的最大高度为多少？
(c) 若小脉冲向下，结果如何？(d) 两脉冲刚过之后，绳子形状如何？
\end{exercise}

\begin{exercise}[$\star\star$]\label{exo:g12:mechanical-waves:8}
(a) 两根晾衣绳受到相同的张力，分别在每根绳上
发出一个脉冲。哪一根脉冲更快：细绳还是粗绳？为什么？(b) 一位
吉他手拧紧一根弦，传播速度会怎样变化？(c) 把抖动
力度加倍（脉冲更高），脉冲会传得更快吗？
\end{exercise}

\begin{exercise}[$\star\star$]\label{exo:g12:mechanical-waves:9}
一名工人在相距 \qty{1.8}{km} 处敲击钢轨，你把一只耳朵贴在
钢轨上听。求声音经钢轨（\qty{5900}{m/s}）
和空气（\qty{340}{m/s}）传来的时间，以及两次响声之间的时间差。
\end{exercise}

\begin{exercise}[$\star\star$]\label{exo:g12:mechanical-waves:10}
一根绳上传播着波长为 \qty{24}{cm} 的周期波。对于下列两点间距，
判断两点是同相、反相还是都不是：(a) \qty{48}{cm}；
(b) \qty{12}{cm}；(c) \qty{36}{cm}；(d) \qty{30}{cm}；(e) \qty{6.0}{cm}？
\end{exercise}

\begin{exercise}[$\star\star$]\label{exo:g12:mechanical-waves:11}
一次地震发出的 P 波速率为 \qty{6.0}{km/s}，S 波速率为
\qty{3.5}{km/s}，两者周期均为 \qty{2.0}{s}。求两种波的波长，
并与一幢建筑作比较，说明为何整个街区会
几乎同步地起伏。
\end{exercise}

\begin{exercise}[$\star\star\star$]\label{exo:g12:mechanical-waves:12}
浮标处水面高度的记录显示，波峰每隔
\qty{0.50}{s} 出现一个；航拍照片显示相邻波峰相距 \qty{0.75}{m}。
(a) 求 $T$ 和 $f$。(b) 求 $\lambda$。(c) 由此推出速率。(d) 浮标
在 $t = 0$ 时恰位于波峰上，下一次何时再位于波峰上？(e) 距浮标 \qty{1.875}{m}
的一点与浮标同相、反相，还是都不是？
\end{exercise}

\begin{exercise}[$\star\star\star$]\label{exo:g12:mechanical-waves:13}
一个锚定的浮标在 \qty{80}{s} 内完成 $10$ 次全振动，
其所处的规则涌浪波峰相距 \qty{120}{m}。(a) 求周期和
频率。(b) 求涌浪速率，分别以 \unit{m/s} 和
\unit{km/h} 表示。(c) 此刻位于浮标处的波峰需多长时间才能到达
相距 \qty{900}{m} 的海滩？(d) 浮标会随之漂上岸吗？
\end{exercise}

\begin{exercise}[$\star\star\star$]\label{exo:g12:mechanical-waves:14}
两个相同的脉冲，一个向上、一个向下，相向传播，
各以 \qty{3.0}{m/s} 的速率运动，中心相距 \qty{6.0}{m}。(a) 两者中心何时
重合，此时绳子呈何形状？(b) 此时绳子是否
静止？能量去了哪里？(c) 接下来会
发生什么？(d) 为什么“瞬时平直”并不意味着“那里什么都没有”？
\end{exercise}

\begin{exercise}[$\star\star\star$]\label{exo:g12:mechanical-waves:15}
一颗小石子激起的涟漪在静水池塘上扩散。(a) 求波峰
在 $r = \qty{0.50}{m}$ 与 $r = \qty{8.0}{m}$ 处的长度；增长倍数为多少？
(b) 波峰的能量近似守恒：每米的能量
与振幅会如何变化？(c) 为什么直线波前的涌浪
衰减得慢得多？(d) 那么体育场“人浪”呢，是否完全不衰减？
\end{exercise}

\section{问题：读懂地震图}

\begin{problem}[{周末问题——读懂地震图：两列波
从断层处向外迸发，三个圆确定震中，一道慢速的
第三波为海岸带来预警}]
\label{pb:g12:mechanical-waves:1}
地震观测站的夜班。09:14:32 时，A 台站的指针
突然跳动：先是尖锐的快速抖动，随后是较慢、幅度较大的摆动。一次
地震在某处使岩石破裂，两列波沿地壳
向外传播：P 波为压缩波，速率 $v_P = \qty{6.0}{km/s}$；S 波为
侧向抖动，速率 $v_S = \qty{3.5}{km/s}$。

\medskip
\noindent\textbf{第一部分——两列波，同一时钟。}
\begin{enumerate}
  \item 根据上述描述，将 P 波和 S 波
        分别归为横波或纵波。
  \item 哪种波先到达每个台站？计算两种波
        在 $d = \qty{120}{km}$ 距离上的传播时间。
  \item 证明 P、S 两波到达的时间差为
        $\Delta t = d\,(1/v_S - 1/v_P)$，并计算
        $d = \qty{120}{km}$ 时的值。
  \item 推导地震学家的测距公式
        $d = \dfrac{v_P\,v_S}{v_P - v_S}\,\Delta t$，并证明该系数
        约为每秒时间差对应 \qty{8.4}{km}。
  \item 为什么时间差随距离增大——这对一个从未看到地震开始的
        台站而言，为何又是一份“馈赠”？
\end{enumerate}

\medskip
\noindent\textbf{第二部分——地图上的圆。}
\begin{enumerate}[resume]
  \item A 台站测得 $\Delta t_A = \qty{25}{s}$，震中
        距 A 台站多远？
  \item 说明为什么这一个读数只能把震中定位在一个圆上，
        而不是某一点。
  \item B、C 两台站测得 $\Delta t_B = \qty{15}{s}$ 与
        $\Delta t_C = \qty{32}{s}$，它们到震中的距离各为多少？
  \item 说明三个圆如何共同确定震中的位置。
  \item P 波于 09:14:32 到达 A 台站。计算它
        从震中传播过来所需的时间，以及地震的发震时刻。
\end{enumerate}

\medskip
\noindent\textbf{第三部分——穿越大洋的波。}
三个圆在近海相交：断层在海床之下破裂，
并使水体开始运动。深水中的长波以
由水深决定的速率传播——在开阔大洋中约为 \qty{200}{m/s}，
在近岸海滩附近约为 \qty{10}{m/s}（在本阶段不加证明地承认，
该公式在本科一年级分册中推出）。某港口位于距震中
$D = \qty{600}{km}$ 处。
\begin{enumerate}[resume]
  \item 将 \qty{200}{m/s} 换算为 \unit{km/h}，并举出一种
        巡航速率与之相当的交通工具。
  \item 海啸需要多长时间才能到达港口？
  \item 港口自己的地震仪先感受到 P 波。经过多久？若警报
        在此时刻拉响，海岸在海啸到来前能获得多少
        预警时间？
  \item 海啸在海上的周期约为 \qty{15}{min}。计算其
        在海上的波长，并说明为何船只察觉不到它。
  \item 接近海滩时，速率降至 \qty{10}{m/s}，而周期
        保持不变。计算新的波长。波在堆积时
        必然发生什么变化？
\end{enumerate}

\medskip
\noindent\textbf{第四部分——能量与报告。}
一次强震释放的能量约为 \qty{1e15}{J}，
由地震波携带。
\begin{enumerate}[resume]
  \item 求距震中 \qty{10}{km} 处圆形波前的长度，
        再求距震中 \qty{210}{km} 处的长度，它增大为原来的多少倍？
  \item 推断即使在无损耗的地壳中，震动仍随距离减弱的原因——
        并指出实际地壳中附加的第二个原因。
  \item 液体不能抵抗剪切。远端台站记录到 P 波，
        但没有 S 波：这说明地球的地核具有怎样的性质？
  \item 上一学年能量一章中的高速列车
        （见\cref{ch:g11:mechanical-energy}）具有约 \qty{1.5e9}{J}
        的动能。\qty{1e15}{J} 相当于多少列这样的列车？
  \item 用两句话撰写报告：地震发生在何处、何时，
        港口得到了多少预警时间？给出三个数值。
\end{enumerate}
\end{problem}
